\section{Related Work}\label{sec:related}

xlog draws on five lines of research---neurosymbolic programming, GPU-accelerated Datalog, probabilistic logic programming, GPU SAT, and differentiable ILP---and is, to our knowledge, the first to make the full symbolic spectrum GPU-resident and uniformly differentiable.

\subsection{Neurosymbolic Programming}

DeepProbLog~\cite{deepproblog} replaces selected probabilistic facts with neural-network outputs, enabling end-to-end gradient training of neural predicates within a probabilistic-logic framework; NeurASP~\cite{neurasp} integrates network outputs as probability annotations on answer-set programs. Both perform symbolic inference on the CPU and shuttle gradients across the bus. Scallop~\cite{scallop} provides a differentiable Datalog with provenance-semiring reasoning~\cite{green2007provenance} and a relational symbolic representation, and is the closest language-level analogue to xlog; its reasoning core, however, is CPU-based. Lobster~\cite{lobster} is the most directly comparable system: it GPU-accelerates neurosymbolic programs in the Scallop style and shares xlog's premise that the symbolic side belongs on the GPU. xlog is complementary along three axes: it covers a broad symbolic spectrum (exact knowledge compilation and epistemic reasoning, not only provenance-based differentiable Datalog); it \emph{verifies} each compiled circuit's logical equivalence to its source formula on-device via CDCL; and it enforces a strict residency contract with byte-level transfer accounting. We make no head-to-head performance claim against Lobster (\Cref{sec:limits}). Differentiable-SAT and logic-as-loss methods---SATNet~\cite{satnet}, semantic loss~\cite{semanticloss}, and Logic Tensor Networks~\cite{ltn}---instead couple neural models to constraints through relaxed or fuzzy surrogates; xlog couples through \emph{exact} weighted model counting, trading their broader applicability for exact gradients and the on-device verifiability of \Cref{sec:prob}.

\subsection{GPU-Accelerated Datalog and Joins}

GPUlog~\cite{gpulog} ported bottom-up fixed-point computation to CUDA using HISA indexing and reported up to $45\times$ speedups over CPU Datalog engines; VFLog~\cite{vflog} advanced this with a vertically-fused, column-oriented kernel design avoiding intermediate materialization, reporting up to $200\times$ over CPU column engines; mnmgDatalog~\cite{mnmgdatalog} extends GPU Datalog to multi-node, multi-GPU clusters. Most relevant to xlog's join path, Sun et al.~\cite{sungpuwcoj} scale worst-case-optimal joins to GPUs. All target deterministic Datalog only: none supports probabilistic semantics, weighted model counting, or gradient computation over derived facts. xlog shares the columnar, kernel-fused execution model and additionally promotes eligible cyclic and multiway rules to a worst-case-optimal join path~\cite{ngo2018wcoj,veldhuizen2014leapfrog,wang2023freejoin}, but treats that engine as one component of a substrate that extends through circuit-based inference to neural-symbolic gradient propagation in a single address space.

\subsection{Probabilistic Logic Programming}

ProbLog2~\cite{problog2} established the modern pipeline for exact probabilistic inference: ground the program, encode provenance as a propositional formula, compile to a tractable target (d-DNNF~\cite{darwiche2001dnnf} or SDD~\cite{sdd}), and evaluate weighted model counts~\cite{chaviraDarwiche2008}. The compilers (D4~\cite{d4compiler}, c2d~\cite{c2d}) run as host-side subprocesses, and recent work makes such compilation \emph{certified}~\cite{capelli2019}. xlog adopts the same compile-once/evaluate-many model but moves the entire pipeline---provenance, Tseitin encoding, D4 compilation, equivalence verification, and forward/backward evaluation---onto the GPU, so probabilistic inference shares one address space with the neural and deterministic layers.

\subsection{GPU SAT}

ParaFROST~\cite{parafrost} parallelizes CDCL clause-database simplification on the GPU while retaining sequential conflict-driven search on the CPU; FastFourierSAT~\cite{fastfouriersat} reformulates SAT as continuous optimization and applies GPU local search, an incomplete solver for satisfiable instances. Both are standalone competition solvers. xlog's GPU CDCL module instead serves as a verification component inside knowledge compilation (\Cref{sec:prob}), proving circuit--formula equivalence via two UNSAT proofs and providing the certificates the probabilistic and epistemic backends rely on.

\subsection{Differentiable ILP}

Differentiable ILP~\cite{dilp} casts first-order rule induction as differentiable optimization, scoring candidate rules by soft forward-chaining; dense materialization over the rule space scales cubically in the number of constants, and implementations remain CPU-based. xlog's dILP subsystem (\Cref{sec:nesy}) replaces dense materialization with sparse GPU bitmasks and on-device scatter--gather credit assignment, avoiding the cubic blowup while preserving differentiability.

\subsection{Logic Programming Languages}

xlog inherits ideas from a long line of logic languages but is the first to target GPU execution of the full surface. Prolog~\cite{swiprolog} offers unification and metaprogramming on a CPU interpreter; Mercury~\cite{mercury} introduced strong typing and modes, compiling to native code; Souffl{\'e}~\cite{souffle} is a typed Datalog compiling to C++ for program analysis; and clingo~\cite{clingo} provides answer-set semantics on CPU grounders and solvers. xlog adopts the typed-predicate discipline but targets device kernels, and unifies typed Datalog, probabilistic facts, neural predicates, and epistemic operators in one language compiled entirely to the GPU.
